\subsection*{Висновки до третього розділу}
\addcontentsline{toc}{subsection}{Висновки до третього розділу}

У третьому розділі дисертаційного дослідження організовано та проведено дослідно-експериментальну роботу з перевірки ефективності розробленої методики навчання розробки імерсивних освітніх ресурсів майбутніх учителів інформатики. Експеримент проводився на базі кафедри інформатики та прикладної математики Криворізького державного педагогічного університету протягом 2020--2025~рр. за участю 64 студентів спеціальності 014 Середня освіта (Інформатика) другого (магістерського) освітнього рівня. Отримано такі основні результати.

\begin{enumerate}[label=\arabic*.]
\item Розроблено критеріальний апарат для оцінювання рівня сформованості вмінь майбутніх учителів інформатики у сфері розробки імерсивних освітніх ресурсів, що включає чотири критерії (мотиваційний, когнітивний, операційно-діяльнісний, рефлексивно-оцінювальний) із відповідними показниками та чотирма рівнями сформованості (низький, середній, достатній, високий). Кожен критерій узгоджено з відповідним блоком структурно-функціональної моделі підготовки (підрозділ~\ref{subsec:2.2}) та педагогічними умовами навчання (підрозділ~\ref{subsec:2.3}), що забезпечує внутрішню цілісність дослідження. Розроблено 100-бальну оцінювальну шкалу з рівномірним розподілом між критеріями (по 25 балів) та комплекс діагностичних інструментів, що включає анкетування, тестування (40 питань), оцінювання лабораторних робіт та проєктів, есе-рефлексії, взаємооцінювання та самооцінювання.

\item Проведено трифазний педагогічний експеримент: констатувальний (2020--2021), формувальний (2021--2024) та контрольний (2024--2025) етапи. На констатувальному етапі підтверджено статистичну еквівалентність експериментальної ($n_1 = 34$) та контрольної ($n_2 = 30$) груп за критерієм $\chi^2$ Пірсона ($\chi^2_{\text{емп}} < \chi^2_{\text{крит}}$ при $p > 0{,}05$). На формувальному етапі в експериментальній групі впроваджено розроблену методику з використанням авторського навчально-методичного забезпечення: навчального посібника <<Розробка WebAR-застосунків для освіти>> (10 розділів), курсу <<Синтетична реальність в освіті>> у LMS Moodle (18 тижневих модулів), лабораторного практикуму та кейсу віртуальної хімічної лабораторії (anchem-ar).

\item Результати контрольного етапу засвідчили статистично значущу позитивну динаміку за всіма критеріями в експериментальній групі порівняно з контрольною. Критерій $\chi^2$ Пірсона для загального розподілу: $\chi^2_{\text{емп}} > \chi^2_{\text{крит}} = 7{,}815$ при $\alpha = 0{,}05$ та $k = 3$; для кожного з чотирьох критеріїв окремо: $\chi^2_{\text{емп}} > \chi^2_{\text{крит}}$. Це дає підстави відхилити нульову гіпотезу про відсутність відмінностей між групами та підтвердити ефективність розробленої методики. Частка студентів ЕГ з достатнім та високим рівнями сформованості вмінь зросла від 23,5\% (констатувальний етап) до 70,6\% (контрольний етап), тоді як у КГ аналогічний показник зріс від 23,4\% до 36,7\%.

\item Найбільш суттєву позитивну динаміку зафіксовано за операційно-діяльнісним критерієм ($\chi^2_{\text{емп}} = 9{,}217$): частка студентів ЕГ з низьким рівнем зменшилася від 41,2\% до 5,9\%, а з достатнім та високим~--зросла від 23,5\% до 70,6\%. Це підтверджує ефективність практико-орієнтованого підходу із систематичним лабораторним практикумом (18 тижнів практичних завдань з WebAR), творчими завданнями та реальними проєктами. Аналіз підсумкових проєктів засвідчив, що 76,5\% студентів ЕГ здатні самостійно розробити WebAR додаток з відстеженням зображень, 64,7\%~--з відстеженням обличчя, 52,9\%~--з відстеженням довкілля, 44,1\%~--з інтеграцією моделей машинного навчання.

\item Апробовано трикрокову методику навчання WebAR з інтеграцією моделей машинного навчання \cite{Semerikov2024WebARML_CoSinE, Semerikov2024WebAR_EduDim}: крок~1~--інтеграція стандартних ML-моделей TensorFlow.js (handpose) для відстеження жестів рук (73,5\% студентів ЕГ успішно виконали); крок~2~--навчання власних моделей за допомогою Google Teachable Machine з експортом до TensorFlow.js (58,8\%); крок~3~--комплексна інтеграція ML-моделей у повнофункціональні WebAR додатки з педагогічним сценарієм (44,1\%). Методика дає змогу поступово формувати у студентів компетентності з ML у контексті імерсивних технологій без необхідності глибоких знань з Data Science.

\item Сформульовано шість груп рекомендацій щодо вдосконалення навчання розробки ІОР: (1)~рекомендації щодо добору технологій (MindAR + Three.js як основна зв'язка; WebXR для відстеження довкілля; A-Frame для прототипування; TensorFlow.js + Teachable Machine для ML); (2)~рекомендації щодо проєктування курсу (послідовність Three.js $\to$ Image Tracking $\to$ Face Tracking $\to$ World Tracking $\to$ ML Integration; обсяг не менше 18 тижнів); (3)~рекомендації щодо оцінювання; (4)~рекомендації щодо інтеграції у Moodle; (5)~рекомендації щодо інтеграції ML; (6)~рекомендації щодо стратегій викладання.

\item Спрямований якісний контент-аналіз 14 транскриптів відеолекцій (1005 сегментів, $\approx$23 години) за шістьма вимірами кодування, побудованими на теоретичних конструктах розділів~1--2, підтвердив реалізацію всіх чотирьох педагогічних умов у навчальному процесі та виявив темпоральну динаміку факторів CAMIL: мотивація домінує на теоретико-аналітичному етапі (26 сегментів), тілесність~--на практико-розвивальному (66 сегментів, пік на тижнях~8--9 при вивченні Face Tracking), саморегуляція концентрується на перехідних тижнях (пік на тижні~10). Когнітивне навантаження управляється рівномірно протягом усього курсу.

\item Аналіз розвитку компонентів TPACK у педагогічному дискурсі виявив домінування TK (344 сегменти), що є закономірним для технологічно-орієнтованого курсу, та обмежену присутність TPK і TPACK-інтегрованого компоненту в лекціях. Водночас кількісні дані (67,6\% студентів ЕГ обґрунтовують педагогічну доцільність ІОР) свідчать, що формування TPK відбувається переважно через практичну діяльність з проєктування навчальних ресурсів, а не через лекційний дискурс.

\item Тріангуляція кількісних та якісних даних (конвергентний паралельний дизайн) продемонструвала збіжність статистичної значущості з якісними свідченнями педагогічних механізмів за всіма критеріями. Зокрема, найвище значення $\chi^2_{\text{емп}} = 9{,}217$ за операційно-діяльнісним критерієм конвергує з домінуванням стратегії моделювання ($\approx$43\%) та піком тілесності при вивченні Face Tracking. Часткова конвергенція за рефлексивно-оцінювальним критерієм пояснюється тим, що саморегуляція формується переважно через письмові форми, не зафіксовані у лекційних транскриптах.

\item Аналіз стратегій викладання виявив домінування моделювання/демонстрації ($\approx$43\%, зокрема live coding) з систематичним зниженням від теоретико-аналітичного (46,9\%) до інтегративно-рефлексивного (38,6\%) етапу. Зростання частки troubleshooting-стратегії (від 14,8\% до 21,3\%) відображає збільшення складності технологій. Загальна тенденція~--від інструктор-центрованих до студент-центрованих стратегій~--відповідає послідовності етапів методики, де роль студента поступово зростає від спостерігача до самостійного розробника.
\end{enumerate}

Результати дослідно-експериментальної роботи підтвердили гіпотезу дослідження: рівень сформованості вмінь майбутніх учителів інформатики у сфері розробки імерсивних освітніх ресурсів суттєво підвищується за умов впровадження розробленої методики навчання, що спирається на структурно-функціональну модель підготовки та реалізацію визначених педагогічних умов. Інтеграція кількісних та якісних методів аналізу дала змогу не лише зафіксувати \textit{що} змінилося (статистична значущість за всіма критеріями), а й розкрити \textit{як} саме навчальний процес забезпечив ці зміни (педагогічні механізми, темпоральна динаміка, стратегії викладання).
